\documentclass[12pt,a4paper]{article}
\usepackage[utf8]{inputenc}
\usepackage[T1]{fontenc}
\usepackage[polish]{babel}
\usepackage{geometry}
\geometry{margin=2.5cm}
\usepackage{enumitem}
\usepackage{titlesec}

% Główne sekcje (np. 1. Postanowienia ogólne) do lewej, czcionka 12pt, pogrubiona
\titleformat{\section}[block]{\raggedright\normalsize\bfseries}{}{0pt}{}

% Paragrafy (np. § 1) wyśrodkowane, czcionka 12pt, pogrubiona
\titleformat{\subsection}[block]{\filcenter\normalsize\bfseries}{}{0pt}{}

\begin{document}

\begin{center}
\textbf{Regulamin praktyk zawodowych dla studiów o profilu praktycznym\\Kierunek Informatyka}
\end{center}
\vspace{1em}

\section*{1. Postanowienia ogólne}

\subsection*{\S 1}
\begin{enumerate}
    \item Regulamin praktyk zawodowych, zwany dalej „Regulaminem", określa zasady organizowania i zaliczenia praktyk zawodowych w Instytucie Informatyki Stosowanej im. Krzysztofa Brzeskiego, zwanym dalej  Instytutem".
    \item Praktyki zawodowe są organizowane zgodnie z Regulaminem Studiów Akademii Nauk Stosowanych w Elblągu.
\end{enumerate}

\subsection*{\S 2}
\begin{enumerate}
    \item Praktyki zawodowe odbywają się w oparciu o Porozumienie w sprawie praktyk zawodowych pomiędzy ANS w Elblągu, zwaną dalej „Uczelnią" a wybranymi w kraju lub/i za granicą jednostkami organizacyjnymi, zwanymi dalej „Zakładami Pracy". Wzór Porozumienia określa załącznik nr 1 do Regulaminu. Dopuszcza się stosowanie druku Porozumienia przedkładanego przez Zakład Pracy.
    \item Wraz z porozumieniem uczelnia wysyła do zakładu pracy Regulamin praktyk zawodowych dla kierunku informatyka oraz Program praktyki (Zał. nr 2.).
    \item Praktyka zawodowa może być realizowana w oparciu o dokumentację zgodną z programem „Erasmus+".
    \item Praktyka zawodowa może być realizowana w jednostce organizacyjnej wskazanej i uzgodnionej przez Uczelnię lub w jednostce organizacyjnej zaproponowanej i uzgodnionej przez studenta - jeżeli spełnia ona warunki zawarte w \S 3 ust. 2.
\end{enumerate}

\subsection*{\S 3}
\begin{enumerate}
    \item Miejsce odbycia praktyki zawodowej musi być uzgodnione z Uczelnianym opiekunem praktyk zawodowych (UOPZ) przed rozpoczęciem praktyki.
    \item Wybór miejsca odbywania praktyki zawodowej i charakter wykonywanej przez studenta pracy musi być zgodny z kierunkiem studiów, aby umożliwić osiągnięcie zakładanych efektów uczenia się i zdobycie umiejętności zawodowych zgodnych z kierunkiem kształcenia.
    \item Opiekunem praktyk ze strony zakładu (ZOPZ) pracy powinien być doświadczony pracownik zatrudniony na stanowisku informatycznym (np. administrator sieci, administrator systemów, programista, referent ds. informatyki itp.) - wymagane wykształcenie wyższe na kierunkach informatycznych lub pokrewnych (automatyka, robotyka, elektronika, telekomunikacja, elektrotechnika, matematyka stosowana); preferowany inżynier informatyk.
\end{enumerate}

\subsection*{\S 4}
\begin{enumerate}
    \item Student, który wykonuje bądź w okresie ostatnich 3 lat wykonywał pracę zawodową, prowadził własną działalność gospodarczą, odbywał staż, a jego zakres obowiązków zawodowych jest/był zgodny z programem praktyki zawodowej, może realizować praktykę zawodową w ramach wykonywanych obowiązków lub wystąpić z wnioskiem o zaliczenie na jej poczet wykonywanych czynności zawodowych.
    \item Student, który odbywać będzie praktykę zawodową w ramach wykonywanej pracy zawodowej bądź występujący z wnioskiem o jej zaliczenie zobligowany jest do przedłożenia odpowiednio: zaświadczenia o zatrudnieniu zakresu obowiązków (opisu stanowiska pracy) wykonywanych w ramach zatrudnienia, prowadzonej działalności gospodarczej, stażu oraz zaświadczenia z CEIDG lub odpisu z KRS albo innych dokumentów, na podstawie których można ocenić nabycie efektów uczenia się przewidzianych dla praktyki zawodowej.
    \item Merytoryczną ocenę wniosku studenta złożonego na podstawie ustępu 2 dokonuje Komisja do zaliczania praktyk zawodowych. Wzór merytorycznej oceny wniosku studenta przez Komisję do spraw praktyk stanowi załącznik 4a. Ostateczną decyzję co do uznania efektów uczenia się podejmuje dyrektor instytutu. Wniosek studenta o uznanie efektów uczenia się zrealizowanych w ramach pracy zawodowej, stażu lub działalności gospodarczej stanowi załącznik Nr 4b.
\end{enumerate}

\subsection*{\S 5}
\begin{enumerate}
    \item Potwierdzenia realizacji zadań wykonywanych w ramach praktyki zawodowej:
    \begin{enumerate}[label=\alph*)]
        \item realizowanej na podstawie porozumienia w sprawie praktyk zawodowych dokonuje Zakładowy opiekun praktyki zawodowej w Dzienniku Praktyk (zał. Nr 6.) i na Karcie praktyki zawodowej, która stanowi załącznik nr 3.
        \item w przypadku realizujących praktykę w ramach pracy zawodowej bezpośredni przełożony w oparciu o dostarczone sprawozdanie załącznik nr 7a,
        \item w przypadku realizujących praktykę na podstawie własnej działalności gospodarczej Uczelniany opiekun praktyk zawodowych w oparciu o dostarczone sprawozdanie załącznik nr 7a.
    \end{enumerate}
    \item Do potwierdzenia efektów uczenia się przewidzianych dla praktyki zawodowej służy załącznik nr 4.
\end{enumerate}

\subsection*{\S 6}
\begin{enumerate}
    \item Egzamin z praktyki zawodowej warunkuje zaliczenie semestru, w którym odbywają się praktyki zawodowe oraz uzyskanie rejestracji na ostatni semestr studiów i dopuszczenie studenta do egzaminu dyplomowego.
    \item W przypadku studenta, który nie realizuje praktyki zawodowej z powodu choroby w czasie łącznym dłuższym niż 40 godzin, ma on obowiązek zwrócić się z wnioskiem o przedłużenie czasu trwania praktyki zawodowej i zrealizować praktykę w wymiarze wnikającym z programu studiów.
    \item Przedłużenie czasu trwania praktyki zawodowej jest możliwe maksymalnie na okres miesiąca po zaplanowanym dniu zakończenia praktyki zawodowej.
\end{enumerate}

\subsection*{\S 7}
Praktykom zawodowym przypisuje się punkty ECTS.

\section*{2. Cele praktyki zawodowej}
\begin{enumerate}
    \item Zapoznanie z działalnością przedsiębiorstw branży informatycznej lub działów informatycznych w przedsiębiorstwach, instytucjach, urzędach, itp.
    \item Zapoznanie z realizacją projektów informatycznych, procesami wdrażania systemów informatycznych lub procedurami związanymi z utrzymaniem urządzeń, obiektów i systemów teleinformatycznych.
    \item Wykonanie samodzielne lub poprzez współudział w zespole minimum jednej pozycji z zakresu:
    \begin{enumerate}[label=\alph*)]
        \item oprogramowania (systemów, aplikacji lub komponentów programowych) lub projektów systemów oprogramowania,
        \item projektów graficznych lub multimedialnych, modelowania 3D, fotogrametrii,
        \item konfiguracji systemów lub projektów sieci zlecanych przez firmy.
    \end{enumerate}
    \item Zdobycie doświadczenia w korzystaniu z norm i standardów stosowanych w informatyce.
    \item Zapoznanie się z zasadami bezpieczeństwa pracy i ergonomii w zawodzie informatyka.
    \item Zdobycie doświadczenia w pracy zespołowej.
\end{enumerate}

\section*{3. Organizacja praktyk zawodowych}
\begin{enumerate}
    \item Praktyka zawodowa realizowana jest na 7. semestrze studiów w wymiarze 6 miesiecy (120 dni roboczych) tzn. 960 godzin. Studentów obowiązuje dzienny wymiar czasu pracy stosowany w danym zakładzie pracy, jednak nie dłużej niż 8 godzin.
    \item Praktykę należy rozpocząć przed semestrem 7. tak, aby zakończyła się do połowy lutego.
    \item W szczególnie uzasadnionych przypadkach Dyrektor Instytutu Informatyki Stosowanej im. K. Brzeskiego, na pisemny wniosek studenta, może wyrazić zgodę na realizację praktyki zawodowej w innym terminie niż wskazany w programie studiów. Zmiana terminu odbywania praktyki zawodowej nie zwalnia studenta z uczestnictwa w zajęciach przewidzianych w programie studiów.
    \item W przypadku gdy osiągnięcie zakładanych efektów uczenia się nie jest możliwe na jednym stanowisku pracy, konieczne jest odbycie praktyki na kilku stanowiskach pracy bądź w różnych zakładach pracy.
    \item Przed przystąpieniem do praktyki, student wraz z Uczelnianym opiekunem praktyk zawodowych (UOPZ) i zakładowym opiekunem praktyk zawodowej (ZOPZ) uzgadniają program i harmonogram praktyki (Zał. 2a).
    \item W przypadku realizacji praktyki w kilku zakładach, konieczne jest prowadzenie odrębnej dokumentacji praktyki.
    \item W przypadku, kiedy student wyjeżdża na studia za granicę w ramach programów unijnych, praktykę zawodową zobowiązany jest zaliczyć we wcześniejszym lub późniejszym terminie. Zgoda na przesunięcie terminu praktyki jest udzielana przez Dyrektora Instytutu Informatyki Stosowanej im.K. Brzeskiego na pisemny wniosek studenta po zasięgnięciu opinii Opiekuna praktyk. Możliwe jest odbywanie praktyk za granicą kraju, pod warunkiem spełnienia wymagań niniejszego regulaminu.
    \item W trakcie trwania praktyki zawodowej przewiduje się minimum jedną hospitację dokonywaną przez Opiekuna praktyk zawodowych w wybranym zakładzie pracy lub przez inne osoby wskazane przez zastępcę dyrektora Instytutu Informatyki Stosowanej im. K. Brzeskiego ds. organizacyjnych i naukowych.
    \item Za organizację praktyk zawodowych odpowiada Opiekun praktyk zawodowych (UOPZ) wyznaczony przez Rektora Uczelni.
\end{enumerate}

\section*{4. Zasady dokumentowania przebiegu praktyki zawodowej oraz procedura jej zaliczania}
\begin{enumerate}
    \item W trakcie praktyki student zobowiązany jest do prowadzenia dziennika praktyki, w którym będzie odnotowywał prace jakie wykonywał w poszczególnych dniach roboczych.
    \item W terminie do 7 dni po zakończeniu praktyki zawodowej, student składa do Opiekuna uczelnianego praktyk zawodowych sprawozdanie wraz z dokumentami określonymi w ust. 4.
    \item W wyjątkowych przypadkach, kiedy student nie ma możliwości realizacji praktyki zawodowej i jej zaliczenia z przyczyn niezależnych od studenta, decyzję o przedłużeniu terminu jej realizacji i zaliczenia podejmuje, na pisemny wniosek studenta i po zasięgnięciu opinii Opiekuna praktyk zawodowych, zastępca dyrektora Instytutu Informatyki Stosowanej im. K. Brzeskiego ds. organizacyjnych i naukowych.
    \item Po zakończeniu praktyki zawodowej Student zobowiązany jest do przedłożenia następujących dokumentów:
    \begin{enumerate}[label=\alph*)]
        \item Dziennika praktyki zawodowej (Zał. nr 6.), w którym będzie odnotowywał prace jakie wykonywał w poszczególnych dniach roboczych z uwzględnieniem zagadnień związanych z osiąganiem efektów uczenia się wymaganych w programie praktyki. Wykonanie poszczególnych zadań, opisanych w dzienniku praktyki, potwierdza zakładowy opiekun praktyki.
        \item Sprawozdania z przebiegu praktyki zawodowej, zawierającego element samooceny w zakresie stopnia osiągnięcia założonych efektów uczenia się, podpisanego przez Zakładowego opiekuna praktyk (Zał. nr 7.).
        \item Zaświadczenia odbycia praktyki zawodowej i opinii z Zakładu Pracy, w którym student odbywał praktykę (Zał. nr 3.),
        \item Potwierdzenia uzyskania efektów uczenia się w ramach zatrudnienia/stażu/praktyki zawodowej (Zał. nr 4.).
        \item Kwestionariusza ankiety oceniający przebieg praktyk zawodowych (Zał. 5).
    \end{enumerate}
    \item Praktyka zawodowa podlega zaliczeniu w formie egzaminu ustnego z oceną.
    \item Egzamin z praktyki zawodowej odbywa się przed Komisją egzaminacyjną powoływaną przez dyrektora Instytutu Informatyki Stosowanej im. K. Brzeskiego. W skład Komisji wchodzą: Przewodniczący, Opiekunowie praktyk. Przewodniczącego Komisji wyznacza dyrektor Instytutu Informatyki Stosowanej im. K. Brzeskiego.
    \item Termin egzaminu z praktyki zawodowej wyznacza przewodniczący Komisji egzaminacyjnych ds. praktyk zawodowych w porozumieniu z Opiekunami praktyk zawodowych.
    \item Podstawą zaliczenia praktyki zawodowej jest protokół sporządzony i podpisany przez Komisję egzaminacyjną ds. praktyk
    \item Protokół sporządzany jest indywidualnie dla każdego studenta i dołączany do akt osobowych studenta (Zał. nr 8.).
    \item Opiekun praktyki zawodowej potwierdza zaliczenie praktyki zawodowej odpowiednim wpisem do systemu komputerowego USOS.
\end{enumerate}

\section*{5. Obowiązki opiekunów i studentów}

\subsection*{\S 8}
\begin{enumerate}[label=\arabic*)]
    \item Do obowiązków Opiekuna praktyk zawodowych należy:
    \begin{enumerate}[label=\alph*)]
        \item opracowanie Programu i harmonogramu praktyki zawodowej (Zał. nr 2a.)
        \item zapoznanie studenta z obowiązującym Programem praktyki zawodowej,
        \item wydawanie skierowań na praktykę zawodową. Skierowanie jest częścią Karty praktyki zawodowej (Zał. nr 3.),
        \item nadzór nad przebiegiem praktyki zawodowej (hospitacje),
        \item prowadzenie dokumentacji praktyk zawodowych,
        \item ocena dokumentacji dostarczonej przez studentów z odbytych praktyk zawodowych,
        \item uznanie efektów uczenia się zakładanych dla praktyki zawodowej (Zał. nr 4.)
    \end{enumerate}
    \item Na wniosek Zakładu Pracy student odbywający praktykę zawodową jest zobowiązany ubezpieczyć się indywidualnie od następstw nieszczęśliwych wypadków na czas jej trwania.
    \item Do obowiązków studenta odbywającego praktykę zawodową należy:
    \begin{enumerate}[label=\alph*)]
        \item zapoznanie się z Regulaminem praktyk zawodowych oraz innymi dokumentami,
        \item w przypadku samodzielnego wyszukania miejsca praktyki na zasadach i w terminie określonym przez uczelnię, student jest zobowiązany dostarczyć do dziekanatu Oświadczenie instytucji (Zał. nr 9.).
        \item aktywne uczestniczenie w praktyce zawodowej i realizacja zadań wyznaczonych przez Opiekuna praktyk,
        \item przestrzeganie dyscypliny pracy oraz obowiązujących w Zakładzie Pracy regulaminów, przepisów BHP i innych,
        \item prowadzenie określonej w Regulaminie dokumentacji praktyk.
    \end{enumerate}
    \item Do obowiązków Komisji egzaminacyjnej ds. praktyk należy przeprowadzenie egzaminu z praktyki zawodowej.
\end{enumerate}

\section*{6. Postanowienia końcowe}
\begin{enumerate}[label=\arabic*)]
    \item W sprawach nieuregulowanych niniejszym Regulaminem decyzje podejmuje dyrektor Instytutu Informatyki Stosowanej im. K. Brzeskiego.
    \item Regulamin praktyk zawodowych obowiązuje studentów rozpoczynających praktykę zawodową od roku ak. 2024/2025
\end{enumerate}

\section*{Załączniki:}
\begin{itemize}[label=--]
    \item Zał. nr 1. Porozumienie-zakład pracy
    \item Zał. nr 2. Program praktyki zawodowej
    \item Zał. nr 2 a Program i harmonogramu praktyki zawodowej
    \item Zał. nr 3. Karta praktyki zawodowej
    \item Zał. nr 4. Potwierdzenie efektów uczenia się
    \item Zał. nr 4a. Potwierdzenie uzyskania efektów uczenia się - merytoryczna ocena wniosku studenta
    \item Zał. nr 4b. Wniosek o zaliczenie efektów uczenia się dla praktyki zawodowej na podstawie pracy zawodowej, stażu lub prowadzenia działalności gospodarczej
    \item Zał. nr 5. Kwestionariusz ankiety
    \item Zał. nr 6. Dziennik praktyki zawodowej
    \item Zał. nr 7. Sprawozdanie z praktyki zawodowej - studia stacjonarne
    \item Zał. nr 7a. Sprawozdanie z praktyki zawodowej - studia niestacjonarne
    \item Zał. nr 8. Protokół egzaminu z praktyki zawodowej
    \item Zał. nr 9. Oświadczenie instytucji w sprawie przyjęcia studenta na praktykę zawodową
\end{itemize}

\end{document}